\subsection{Einfluss von Induktivität und Kapazität}

Bei Gleichrichtern werden Induktivität L und Kapazität C eingesetzt, um:
\begin{itemize}
\item die Welligkeit der Gleichspannung zu reduzieren,
\item die Stromform zu beeinflussen,
\item Netzrückwirkungen und Harmonische zu verändern.
\end{itemize}

Man unterscheidet grundsätzlich:
\begin{itemize}
\item \textbf{C-Glättung (kapazitive Last)},
\item \textbf{L-Glättung (induktive Last)},
\item \textbf{LC-Glättung (Filter)}.
\end{itemize}

\subsubsection{C-Glättung (Kondensator am Ausgang)}

Grundprinzip:
\begin{itemize}
\item Der Kondensator lädt sich auf das Spannungsmaximum auf.
\item Zwischen den Leitphasen entlädt er sich über die Last.
\item Die Diode leitet nur kurzzeitig um das Maximum.
\end{itemize}

Typische Eigenschaften:
\begin{itemize}
\item Ausgangsspannung näher am Scheitelwert $\hat U$.
\item Kleiner Spannungsrippel bei grossem C.
\item Sehr kurze, hohe Ladeimpulse im Diodenstrom.
\end{itemize}

Näherung für den Spannungsrippel:
\[
\boxed{\Delta u \approx \frac{I_d}{f C}}
\]

Mittelwert der Gleichspannung:
\[
\boxed{u_d \approx \hat U - \frac{\Delta u}{2}}
\]

Konsequenzen:
\begin{itemize}
\item Hoher Spitzenstrom in den Dioden.
\item Starke Netzstromverzerrung.
\item Schlechter Leistungsfaktor.
\end{itemize}

\subsubsection{L-Glättung (Drossel in Serie)}

Grundprinzip:
\begin{itemize}
\item Die Induktivität erzwingt nahezu konstanten Gleichstrom.
\item Der Strom fliesst kontinuierlich über die ganze Periode.
\item Die Spannung passt sich dem Strom an.
\end{itemize}

Typische Eigenschaften:
\begin{itemize}
\item Gleichstrom nahezu konstant.
\item Spannung stark gewellt.
\item Netzströme rechteckförmig (bei B6: 120°-Leitung).
\end{itemize}

Im stationären Betrieb gilt:
\[
\boxed{\overline{u_L} = 0}
\]

Daraus folgt für ideale Gleichrichter:
\[
\boxed{u_d = \text{Mittelwert der gleichgerichteten Spannung}}
\]

Konsequenzen:
\begin{itemize}
\item Geringerer Dioden-Spitzenstrom.
\item Besserer Leistungsfaktor als bei C-Glättung.
\item Typischer Industrie-Betrieb.
\end{itemize}

\subsubsection{LC-Glättung (Tiefpassfilter)}

Kombination aus:
\begin{itemize}
\item Serieninduktivität L,
\item Parallel-Kondensator C.
\end{itemize}

Wirkung:
\begin{itemize}
\item L glättet den Strom.
\item C glättet die Spannung.
\item Stark reduzierte Restwelligkeit.
\end{itemize}

Grenzfrequenz des LC-Filters:
\[
\boxed{\omega_0 = \frac{1}{\sqrt{LC}}}
\]

Dämpfung hoher Harmonischer:
\[
|H(j\omega)| \downarrow \quad \text{für} \quad \omega \gg \omega_0
\]

\subsubsection{Vergleich C-Glättung vs. L-Glättung}

\begin{center}
\begin{tabular}{l|c|c}
 & C-Glättung & L-Glättung \\
\hline
Stromform & Impulsförmig & Nahe konstant \\
Spannungsrippel & klein & gross \\
Diodenbelastung & hoch & gering \\
Leistungsfaktor & schlecht & besser \\
Industrieeinsatz & Netzteile klein & Leistungsantriebe \\
\end{tabular}
\end{center}

\subsubsection{Typische Prüfungsfälle}

Häufige Aufgabenstellungen:
\begin{itemize}
\item M1U oder B2U mit grossem C $\Rightarrow$ kurze Leitwinkel, Spitzenstrom berechnen.
\item B6U mit grossem L $\Rightarrow$ konstanter Gleichstrom, 120°-Netzstrom.
\item Berechnung von $\Delta u$ aus $I_d, f, C$.
\item Vergleich der Netzstromharmonischen bei C- und L-Glättung.
\end{itemize}

Merksätze:
\begin{itemize}
\item C glättet Spannung, verschlechtert Netzstrom.
\item L glättet Strom, verbessert Netzrückwirkung.
\item Industrie: fast immer L-Glättung.
\end{itemize}
